\section{Declared Data}\label{sec:data}

A contract reads two things and nothing else: declared data and stamped
observations. This section is about the first. Declared data is the part of the
log that is authored deliberately, ahead of the events that consume it --- the
terms of an instrument, the way each of its fields reacts to a corporate action,
the kinds of observation the system will accept, and the operators that carry an
old number across a boundary into a new frame. The discipline throughout is the
same one the door already showed: fix the answer in advance, as data, so that no
event has to improvise.

\subsection{Terms, and versions that only ever append}

Every unit has \emph{ProductTerms}: the terms of the instrument, recorded as a
non-empty, append-only list of versions. The first version is written when the
unit is registered; a unit that is registered but has no terms is not a state the
types allow. There is no operation that edits a terms version, and --- this is the
part that matters for corporate actions --- an adjustment never edits terms
either. When a split halves a strike, it does not overwrite the strike field; it
\emph{appends} a new terms version carrying the halved strike, tagged with the
boundary that caused it. The old version stays exactly as it was written. Terms
are declared data, and declared data, like everything in the log, grows only by
appending.

\subsection{Four dimensions, and how a field reacts}

Why does a split halve a strike but not a coupon, and halve a cash-per-share
dividend but not a percentage one? Because a field's reaction to a corporate
action is fixed entirely by \emph{what kind of quantity the field is}, and there
are only four kinds. Every terms field declares, at the point it is defined,
exactly one \emph{dimension}:

\begin{itemize}[leftmargin=1.5em]
\item a \textbf{quantity} of the unit --- a share count, a number of contracts;
\item a \textbf{price} in the unit's basis --- a strike, a stored reference spot;
\item a \textbf{cash} amount --- a fixed coupon, a notional;
\item a \textbf{dimensionless} value --- a ratio, a percentage, a participation.
\end{itemize}

The dimension is the field's constructor: a field with no dimension is not
invalid, it is unconstructible. A corporate action declares a quantity factor
$f \neq 0$ and a per-share cash amount $c$, and each dimension has one default action,
the same for every product:

\[
  \text{quantity } q \mapsto q \cdot f, \qquad
  \text{price } p \mapsto (p - c)/f, \qquad
  \text{cash and dimensionless: unchanged.}
\]

The factor is nonzero, so the boundary is invertible: a datum carried across it
can be carried back. That is the whole of the constraint, and it is why $f = 0$ is
refused --- a zero factor is non-invertible, collapsing every quantity to zero with
no way to recover the count it came from, so a datum read across such a boundary
could never be transported in either direction. A negative factor is well-formed;
only zero is refused.

Check the arithmetic against intuition. A two-for-one split is $f = 2$, $c = 0$: a
holding of 1{,}000 becomes 2{,}000, a strike of \EUR{100} becomes \EUR{50}, a fixed
coupon in cash is untouched. The price action earns its keep on market data at an
ordinary dividend. An ordinary \EUR{2} dividend is $f = 1$, $c = 2$; a market price
printed at \EUR{102} the instant before the stock goes ex, read across the
boundary, becomes $(102 - 2)/1 = \EUR{100}$ --- exactly the print going
ex-dividend. That is the map doing its one job on a datum in the price dimension.

A price that is a \emph{terms field} --- a stored strike or reference spot on the
instrument itself --- is a different object, and under an ordinary dividend it does
not move: it stands at \EUR{102}. The arithmetic is the same; what differs is that
applying it to a terms field here would pay the same cash twice, and the next
subsection makes that precise. The dimensional action is where a field starts; the
schedule, developed there, is what holds a price field fixed exactly where moving
it would double-count.

The reason one rule suffices for every product is a small theorem about these four
dimensions, and it is the load-bearing fact of the whole adjustment machinery.

\begin{theorem}[Dimension invariance]\label{thm:dim-invariance}
For a quantity $q$ and a price $p$ crossing a boundary $(f, c)$,
\[
  (q \cdot f)\cdot\frac{p - c}{f} \;=\; q\cdot p - q\cdot c .
\]
The marked value after the boundary equals the value before, less exactly the cash
the same boundary paid out.
\end{theorem}

The identity holds field by field, for every product, at every boundary on the
menu, because it is quantified over the four dimensions once and not over products
or event types. A product registered next year adds fields with dimensions; it
adds no cases to this theorem. That is the sense in which the declarations
\emph{relocate} the expert judgement rather than removing it: someone must still
say, correctly, which dimension each field of each product has --- but having said
it once, they never write another adjustment rule, and $N$ products crossing $M$
kinds of corporate action need $N$ sets of field declarations, not $N \times M$
hand-written cases.

\subsection{The adjustment schedule and its three layers}

Declaring dimensions gives every field a \emph{default} reaction. But real term
sheets deviate. A specific note may say its strike adjusts on a special dividend
only above a threshold; a specific determination may be left to a calculation
agent to fix later. So every unit's terms carry, from the first version, an
\emph{adjustment schedule}: for every class of corporate action and every terms
field, exactly one entry saying how that field reacts. An entry comes from exactly
one of three layers, most specific winning:

\begin{enumerate}[leftmargin=1.5em]
\item the \textbf{class default} --- fixed once, in this specification, per class
and dimension: quantities scale by $f$, cash and dimensionless fields stand, and a
price terms field moves to $(p-c)/f$ under every class \emph{except} an ordinary
dividend, where it stands. That one exception is neither market custom imported nor a
consequence of Theorem~\ref{thm:dim-invariance} alone. The re-expression
$p \mapsto (p-c)/f$ preserves value only when the same transaction pays the holder of
the shifted field the matching cash leg $q \cdot c$; a market quote crossing the seam
has that pairing, because the shareholder is paid, but a stored strike has no holder
being paid, so the theorem by itself licenses no strike shift in either class. What
decides the default is what the strike already contains. An ordinary, anticipated
dividend sits in the forward the strike was struck against, so the strike stands ---
shifting it would create option value equal to the cash the boundary has already paid
the shareholder, one distribution counted twice in the system's value accounting. A
special dividend sits in no forward, so a standing strike would impose exactly that
unpriced transfer in the opposite direction, and the strike moves to $(p-c)/f$ to hold
the contracted position invariant. One principle yields both: a price terms field
shifts by precisely the part of the distribution not already embedded in the reference
it was struck against --- zero for an ordinary dividend, all of $c$ for every other
class. The rule is read off that principle, not copied from a convention.
\item a \textbf{declared override} --- a rule in a closed, first-order language
that reads the boundary's parameters and the field's value and evaluates, with no
recursion and no partial operations, to exactly one value. This is where a
bespoke term sheet writes down its deviation.
\item \textbf{designated discretion} --- the schedule names who owes a
determination; until that party's attested \emph{resolved terms} arrive
referencing the boundary, the affected fields have no adjusted value, the boundary
is pending for them, and any attempt to consume them is quarantined.
\end{enumerate}

Adjusted terms enter the log from these two sources --- the schedule's evaluator
applying a declared entry, and the resolved terms of a designated party --- and
from nowhere else.

What makes this safe rather than merely organised is a totality gate. A
registration, and every later terms append, is admitted only if the schedule
yields exactly one well-formed entry for every pair of (corporate-action class,
field) in the version being written; an append that introduces a new field must
extend the schedule in the same transaction, or it is refused. The set of
corporate-action classes is closed and finite --- eight members: split, ordinary
dividend, special dividend, spin-off, stock merger, elective merger, termination,
and rebalance --- so the gate is a finite check, and it is strict. An ill-formed
override is refused, never quietly defaulted. The consequence is a strong one:
past the moment a unit is registered, ``the framework does not know how to handle
this event for this product'' is not a state the system can reach. Either every
class $\times$ field pair has an answer, or the version was never admitted.

The eight are closed by design, and closed does not mean incomplete: an event that
looks like a ninth usually decomposes onto a member. A rights issue is the standing
example. It is a spin-off --- the right is a tradable series distributed to holders
and registered as the child unit --- not a new class of its own. A notice that
genuinely falls outside the eight is refused, fail-closed; the taxonomy is widened
only by changing this specification, never by an incoming datum.

A word on where the schedule is evaluated, because it decides what the door can and
cannot catch. The schedule's language is one definition, fixed once here, but it is
executed at three distinct sites. The door executes it to check \emph{totality} ---
that exactly one well-formed entry exists for every class $\times$ field pair --- and
nothing more; it never evaluates an entry down to a resolved value. The
corporate-action processor executes it to \emph{produce} the resolved terms at a
boundary. And an independent recomputation, run later against the committed log,
executes it a third time to \emph{verify} those terms. The door's evaluation proves
an answer exists and is unique; it does not compute the answer, and it does not
check the number the processor appended. That is precisely why a processor that
appends a wrong resolved value --- a perfectly well-formed terms version carrying
the wrong strike --- passes the door and is caught only afterward, by the
recomputation re-running the same evaluation and comparing.

\subsection{The datum-kind registry}

Observations from the outside --- a settlement price, a dividend forecast, a
barrier fixing --- enter the log as \emph{stamped observations}. But not any
observation may enter. Each \emph{kind} of consumable datum is registered once,
before any datum of that kind can be recorded, and the registration fixes the
kind's identity and declares a dimension for each of its components. A datum is an
observation, never a holding, so its components draw on three dimensions --- price,
cash, dimensionless --- but never quantity.

The reason a datum is registered per component, rather than named as a single
number, is that one datum's components can have different fates under the same
operator. A dividend forecast is a good example. Suppose a forecast lists two
components. The first is an amount of \EUR{3.00} per share; being quoted in the share's price
dimension, it takes the price action under a split, and it is declared in the
\emph{price} dimension. The second is a proportional part, $1\%$ of price, declared
\emph{dimensionless}. A two-for-one split applies $f = 2$ to the whole datum, but
the two components react by their own dimensions: the per-share amount halves to
\EUR{1.50}, the $1\%$ stands. And both contributions are invariant, which is the
point ---
$1{,}000 \times 3 = 2{,}000 \times 1.50 = \EUR{3{,}000}$ on the cash component, and
$1\% \times 100 \times 1{,}000 = 1\% \times 50 \times 2{,}000 = \EUR{1{,}000}$ on
the proportional one. A forecast settles nothing; it feeds a valuation. Had the
proportional component been mis-declared as a price, the split would have halved it to
$0.5\%$ and contributed \EUR{500} to the mark where \EUR{1{,}000} is correct ---
half the datum's contribution, from one wrong word in a declaration. This is
exactly why each registry entry must carry a mandatory invariance witness, and it
is worth saying precisely what the witness checks. It takes each component's
declared dimension and a scale boundary $(f, c)$ and tests the value-invariance
identity for that dimension: the quantity the component delivers must be the same
before and after the boundary once the component is carried across by the map its
declared dimension fixes. A component declared per-share in the price dimension
satisfies $1{,}000 \times p = 2{,}000 \times (p/2)$ under a two-for-one; a component
declared as an absolute cash amount is left unscaled and satisfies the identity only
if it truly is a flat amount that does not depend on the share count. A per-share
amount wrongly declared as absolute cash --- or a flat amount wrongly declared as
per-share --- fails its identity under the replayed boundary, and the entry is
refused at registration. The declaration is a trust assumption, and the witness
bounds that trust to precisely its own width: a dimension that would misbehave under
a scale cannot be registered in the first place.

A kind that was never registered is unconsumable: the door refuses it outright, as
an unregistered kind, returning the payload with the refusal but admitting nothing. This is distinct
from a registered kind that a particular boundary's declaration simply does not
mention --- that datum enters the log and is stopped only when it tries to cross
that boundary. A registry entry, like a terms version, is immutable; a wrong
declaration is corrected by registering a new kind and retiring the old one to new
ingestion, never by editing the entry in place.

\subsection{Operators: one declaration, two consumers}

The last piece of declared data is the operator that carries a number across a
boundary. When a corporate action is recorded, its declaration is a finite map
from datum kinds to \emph{operators}, drawn from a closed menu:

\begin{lstlisting}
data OpSpec
  = Scale     Rational  -- multiply a quantity/price by a nonzero factor
  | Shift     Rational  -- subtract a per-share cash amount from a price
  | Subst     UnitId Rational -- succeed into another unit at a declared ratio
  | Recompose Basket    -- resolve onto a basket of successor units
  | AId                 -- identity, declared -- never assumed
  | Pending             -- boundary known, parameter not yet published
  | Terminal            -- the series ends; no conversion either way
\end{lstlisting}

The menu is closed: a declaration naming anything outside it is refused at the
door, and widening the menu is a change to this specification, not something data
can do. The menu is closed as an alphabet, not as a grammar: a declaration may
bind a datum kind to a composite of menu operators --- a succession at a declared
ratio with a per-share cash leg is $\texttt{Subst}(Q, f) \circ \texttt{Shift}(-c)$,
the interpretation of the single declared pair $(f, c)$ as the price action
$p \mapsto (p - c)/f$ --- and closure means no operator outside the menu, never that
members cannot compose. Two features of it repay attention. First, identity is a declared member,
\texttt{AId}, and is never assumed. A declaration that omits a datum kind does not
pass that kind through unchanged; it provides no conversion for it, and the door
fails closed. The two-case function is exactly this:

\begin{lstlisting}
convert :: Maybe OpSpec -> Datum -> Either Refusal Datum
convert Nothing   _ = Left NoDeclaration   -- omission is refusal
convert (Just op) d = apply op d           -- AId is a `Just`, and is identity
\end{lstlisting}

The difference between \texttt{Nothing} and \texttt{Just AId} is the difference
between ``no answer, stop'' and ``the declared answer is: do nothing''. Only the
second lets a datum cross. Second, the declaration is an operator, not a name. A
name like ``two-for-one split'' would have to be interpreted, identically, by
every consumer forever; an operator like \texttt{Scale 2} is simply applied. This
is what lets one declaration serve two consumers.

A boundary's declared operators are, in fact, consumed twice. They are consumed
once at the boundary's effective time, by the corporate-action contract that
produces the boundary's one-time transaction; and then they stand forever after
at the read seam, converting every pre-boundary datum that is later read across
the boundary. An example makes the two uses concrete. This read-seam application is
the one generic computation the thesis of Section~\ref{sec:thesis} named. Consider a
single day. The log already holds May's \EUR{2} ordinary dividend, recorded as
\texttt{Shift} $-2$. At 09:59 a vendor prints a price of 100, stamped in the basis
current at that moment --- the pre-split frame, which the 10:00 split will leave one
boundary back. At 10:00 a two-for-one split takes effect, and the
corporate-action processor consumes the declaration \texttt{Scale 2} \emph{once}:
it moves 1{,}000 shares to 2{,}000, sets the new basis, appends the strike from
\EUR{100} to \EUR{50} and the stored reference spot from \EUR{100} to \EUR{50},
each tagged with the boundary. At 10:03 someone reads the 09:59 print in the new
frame, and now the same \texttt{Scale 2} is consumed a \emph{second} time, at the
read seam: the stale 100 is converted to $100/2 = 50$, and only then is it allowed
to meet the 2{,}000-share count. The ledger did not price anything and did not
interpret the split; it applied the declared operator, and that is the whole of
the generic work it does.

The composition of operators across several boundaries is a computation, not a
choice. A May print of 102 read across both the dividend and the split composes in
effective order: $(102 - 2)/2 = 50$. The reverse order, $102/2 - 2 = 49$, is not a
valid conversion: no pair of bases stands in that relation, and the declared data
offers no way to assemble it, so the arrival order of a datum can never change the
number it converts to.
Between any two bases on one chain there is at most one conversion per datum kind,
and it exists exactly where every intervening declared conversion does. The
\EUR{50} is not the ledger's opinion. It is the only number the declared data
admits.
